\paragraph{Supplement to Finding A51: fields of quintic iteration and available canonical detectors.}
Let $k=\mathbf F_q$ have characteristic two, let
$\Sigma=\{2,3,5,13\}$, and set
\[
 k_\Sigma=\bigcup_{\substack{m\ge1\;:\text{ every prime}\,
                              \text{dividing }m\text{ lies in }\Sigma}}
                              \mathbf F_{q^m}.
 \tag{FI.1}
\]
The pointed geometric quintic-cover and reflection choices of A33,
iterated from a genus-two curve and a rational point over $k$, have
representatives over $k_\Sigma$. Nevertheless, on every fixed
genus-two target defined over a finite subfield of $k_\Sigma$,
there are infinitely many points in this same field extension whose
canonical class has nonzero $3$-primary component. These two assertions
do not force the prescribed iterated reflection images to meet those
points.

\begin{proof}
\emph{1. One uniform degree support for the geometric choices.}
Let $(C,P)$ be defined over a finite field $K\subset k_\Sigma$.
Its principal polarization gives a nondegenerate alternating Weil
pairing on $J_C[5]$, and Frobenius acts as a similitude, with
multiplier $\#K$ modulo $5$. Consequently all $5$-torsion becomes
rational over an extension of degree dividing
\[
 \begin{aligned}
 G_5&=\#\operatorname{GSp}_4(\mathbf F_5)
      =(5-1)5^4(5^2-1)(5^4-1)\\
    &=37{,}440{,}000=2^9\,3^2\,5^4\,13.
 \end{aligned}
 \tag{FI.2}
\]
For completeness, symplectic bases in dimension four over
$\mathbf F_t$ are counted by choosing a nonzero first vector
in $t^4-1$ ways, its pairing-one partner in $t^3$ ways, then
a nonzero vector in their two-dimensional orthogonal complement
in $t^2-1$ ways and its partner in $t$ ways. This gives
$\#\operatorname{Sp}_4(\mathbf F_t)=t^4(t^2-1)(t^4-1)$;
the surjective multiplier adds the factor $t-1$.
Lagrange's theorem gives the divisibility of the Frobenius order.

Over that extension every geometric cyclic quotient of degree five is
Frobenius invariant with constant deck group. The section at $P$
extends the geometric quotient of the abelianized fundamental group
to an arithmetic quotient by sending that section to zero. This
constructs its constant, $P$-split model; all five lifts of $P$
are rational. Equivalently, among the constant cyclic twists the
$P$-split choice is the unique one, as in A48. Thus specifying a
geometric cover and a geometric lift of $P$ adds no obstruction
to representing the pointed isomorphism class in this field.

The hyperelliptic involution of $C$ is defined over $K$ and acts
as inversion on the cyclic quotient. It has five lifts to the
cover, all involutions, as in A33. Write one as $\tau$ and a deck
generator as $\sigma$. Because the deck group is now constant,
Frobenius acts on the lift set by
\[
                  \tau\longmapsto\sigma^b\tau
                  \qquad(b\in\mathbf Z/5\mathbf Z).
 \tag{FI.3}
\]
This is a translation of that five-element torsor, so an extension
of degree dividing five defines every reflection. Their genus-two
quotients and all images of the five rational lifts are defined
there too. The local quotient/genus argument is the one already
proved in A33.

Thus one step can be represented after a degree dividing
\[
                  B=5G_5=187{,}200{,}000.
 \tag{FI.4}
\]
Inductively every depth-$j$ choice has a representative over
$\mathbf F_{q^{B^j}}$: every intermediate relative degree divides
$B$, and the same field contains all such relative subextensions.
This asserts a common prime support, not a single finite field
for all depths. It concerns geometric pointed choices represented
by their split models, not arbitrary preselected equations or
arbitrary forms over unrelated fields. If another constant twist
is first specified, its five-element fibre torsor is killed by
an extension of degree dividing five, which has the same prime
support.

\emph{2. Canonical $3$-primary detectors already occur densely inside
that constant extension.}
Fix any genus-two target $E/K$ with finite $K\subset k_\Sigma$.
An extension inside $k_\Sigma$ rationalizes $J_E[3]$, since
\[
 \#\operatorname{GSp}_4(\mathbf F_3)
       =103{,}680=2^8\,3^4\,5.
 \tag{FI.5}
\]
After a further extension of degree a power of five, the Weil
bound gives a rational point $O\in E(K)$; enlarge $K$ accordingly.
Put $Q=\#K$, $G=J_E(K)[3^\infty]$, and $d=\#G\ge3^4$.
If $\Phi$ is its Frobenius, it acts as the identity on $J_E[3]$.
For $n=5^m$ the polynomial
\[
  V_n(\Phi)=1+\Phi+\cdots+\Phi^{n-1},\qquad
  \Phi^n-1=(\Phi-1)V_n(\Phi)
 \tag{FI.6}
\]
is invertible on $3$-power torsion: modulo three it is $nI$,
and a nonzero kernel at a higher finite level would contain an
element of order three. It follows that, as actual subgroups,
\[
                  J_E(\mathbf F_{Q^{5^m}})[3^\infty]=G.
 \tag{FI.7}
\]
Apply the already proved fixed Lang-cover count A42 with the
singleton prime set $\{3\}$ and the Abel map $R\mapsto[R-O]$.
Here its modulus is $3$, so every $n=5^m$ is admissible. For
every $h\in G$ it gives
\[
 \left|\#\{R\in E(\mathbf F_{Q^n}):[R-O]_3=h\}
                 -\frac{Q^n+1}{d}\right|
       \le\frac{2(d+1)}d Q^{n/2}.
 \tag{FI.8}
\]
This uses the connected Lang quotient of degree $d$ and genus
$d+1$, with its exact norm scaling; no unproved density
principle is substituted for A42.

For the canonical class $\kappa_R=[\omega_E(-2R)]$ one has
\[
             (\kappa_R)_3=(\kappa_O)_3-2[R-O]_3.
 \tag{FI.9}
\]
Exactly one $h\in G$ makes the right side zero, since
multiplication by two is invertible on $G$. Each of the other
$d-1$ components occurs on unboundedly many points by (FI.8).
All these points belong to $E(k_\Sigma)$ and have a nonzero
canonical $3$-primary component. Their relative density along
these fields tends to $(d-1)/d$, and in particular they are
Zariski dense. Since $\operatorname{ord}_3(2)=2$, A29 applies
to each such pointed target and supplies its prime-power
canonical detector.
\end{proof}

The pro-two or single-auxiliary-prime constant-tower restrictions
in A44 concern a fixed Jacobian and its Frobenius polynomial;
they do not uniformly restrict this changing family of reflection
Jacobians. The preceding calculation gives the sharper relevant
endpoint: the common field support is compatible with abundant
good canonical points on every fixed target. It still does not
force any prescribed finite set of reflection images, or the
full iterated orbit from a specified point, to contain one.
Neither a Hecke-orbit theorem nor the Lang count establishes
that missing incidence. No main-gate reduction is claimed.
